\handbooktitle{EF Core Plugin-Entwicklung · Referenz und Praxis}{Das EF Core Plugin-Handbuch}
  {Von der Options-Extension bis zum NuGet-Paket: alles, was man für professionelle und auch
  sehr komplexe Erweiterungen von Entity Framework Core braucht. Durchgängig erklärt an
  einem vollständigen Plugin für Mandantenfähigkeit mit PostgreSQL und SQL Server.}
\begin{handbookmeta}
  Ausgabe & \handbookversion \\
  Versionen & EF Core 10, Npgsql.EntityFrameworkCore.PostgreSQL 10, .NET 10 \\
  Beispiel & Vollständige Solution unter
  \textbf{github.com/guenni81/EF-Core-Plugin-Handbook}: Kernpaket, zwei Provider-Pakete,
  zwei Beispiel-Apps, Unit- und Integrationstests \\
  Geprüft & Der Beispielcode wird gegen EF Core 10 und Npgsql 10 gebaut; Unit- und
  Integrationstests laufen gegen echtes PostgreSQL und SQL Server. \\
  Aufbau & 9 Teile, 42 Kapitel, 4 Anhänge \\
  Lizenz & MIT-Lizenz, © 2026 guenni81. Frei nutzbar, kopierbar, veränderbar und
  weitergebbar, auch kommerziell; siehe letzte Seite. \\
\end{handbookmeta}

\section*{Über dieses Buch}

Viele Erweiterungen für EF Core beginnen als ein paar Extension-Methods und wachsen dann
unkontrolliert. Wer verstanden hat, wie EF Core intern aufgebaut ist, kann dagegen
Erweiterungen schreiben, die sich wie ein Teil von EF Core anfühlen: Sie werden über die
Options konfiguriert, laufen im internen Service Provider, übersetzen eigene Funktionen nach
SQL, erzeugen Migrationen, schreiben lesbare Snapshots und lassen sich mit anderen Plugins
kombinieren.

Dieses Buch erklärt jeden dieser Bausteine in zwei Schichten. Zuerst kommt das
\textbf{Prinzip}: Wie funktioniert der Mechanismus in EF Core, welche Varianten gibt es,
worauf muss man achten? Danach folgt die \textbf{Umsetzung} im Beispiel-Plugin, mit dem
Code, der auch im Repository steht. Die Kapitel bauen aufeinander auf, lassen sich aber auch
einzeln als Nachschlagewerk lesen. Die Codeauszüge dieser deutschen Fassung haben deutsche
Kommentare; im Repository sind Kommentare und Meldungen englisch, der Code ist sonst
identisch. Das Label jedes Auszugs nennt die Datei.

\begin{datatable}{H{0.50}H{1.50}}
  \textbf{Kasten} & \textbf{Bedeutung} \\
  \tablehead
  \textcolor[HTML]{D99A00}{\textbf{Wichtig}} & Fehlerquellen, die zu falschem Verhalten oder
  Sicherheitslücken führen. \\
  \tablerule
  \textcolor[HTML]{2E8B57}{\textbf{Tipp}} & Bewährte Praxis, die Arbeit spart oder die
  Qualität erhöht. \\
  \tablerule
  \textcolor[HTML]{1D4E89}{\textbf{Hinweis}} & Hintergrund, Einordnung und Verweise auf
  andere Kapitel. \\
\end{datatable}

\begin{note}[Öffentliche und interne APIs]
  Einige APIs, die Plugins brauchen, liegen in Namespaces mit \code{.\allowbreak{}Internal}.
  EF Core nennt sie „pubternal“: öffentlich sichtbar, aber ohne Kompatibilitätsgarantie. Das
  Buch kennzeichnet diese Stellen und erklärt in
  Kapitel~\ref{ch:versioning-and-compatibility}, wie man damit umgeht.
\end{note}

\tableofcontents

% ============================================================
\setpartpreamble{\partintro
  {Was ein Plugin ist, wie EF Core intern aufgebaut ist und welches Szenario uns durch das
  Buch begleitet.}
  {
    \item Kapitel~\ref{ch:what-is-an-ef-core}: Was ist ein EF-Core-Plugin?
    \item Kapitel~\ref{ch:ef-core-s-architecture-from}: Die Architektur von EF Core aus
    Plugin-Sicht
    \item Kapitel~\ref{ch:the-sample-scenario}: Das Beispielszenario
  }}
\part{Grundlagen}\label{part:I}

% ------------------------------------------------------------
\chapter{Was ist ein EF-Core-Plugin?}\label{ch:what-is-an-ef-core}

Der Begriff „Plugin“ wird für sehr unterschiedliche Dinge verwendet. Für dieses Buch hilft
eine Einteilung in vier Stufen, die sich darin unterscheiden, wie tief sie in EF Core
eingreifen:

\begin{datatable}{H{0.65}H{1.41}H{0.94}}
  \textbf{Stufe} & \textbf{Beschreibung} & \textbf{Beispiele} \\
  \tablehead
  1 · Hilfsmethoden & Extension-Methods auf \code{Db\allowbreak{}Context} oder
  \code{IQueryable}, die nur öffentliche API aufrufen. & Paging-Helfer,
  \code{Where\allowbreak{}If}, Seed-Methoden \\
  \tablerule
  2 · Modell-Helfer & Fluent-API-Erweiterungen und Annotations, ohne eigene Dienste in EF
  Core. & Eigene \code{Has\allowbreak{}Xyz(\allowbreak{})}-Methoden, ModelBuilder-Helfer \\
  \tablerule
  3 · Echte Plugins & Eigene Dienste im internen Service Provider: Conventions,
  Query-Übersetzung, Type Mappings, Migrations, Design-Time. & NetTopologySuite-Plugin,
  NodaTime-Plugin, EFCore.NamingConventions \\
  \tablerule
  4 · Provider & Komplette Datenbankanbindung inklusive Verbindung, Update-Pipeline und
  SQL-Generierung. & Npgsql, Pomelo (MySQL), SQLite, Cosmos \\
\end{datatable}

Dieses Buch behandelt \textbf{Stufe 3} vollständig. Stufen 1 und 2 brauchen kein eigenes
Buch: Sie nutzen nur die öffentliche API von EF, und
Kapitel~\ref{ch:annotations-and-fluent-api} behandelt, was an Fluent-API-Erweiterungen
speziell ist. Die Werkzeuge von Stufe 4 (etwa der Austausch des SQL-Generators) werden dort
erklärt, wo Plugins sie brauchen.

\section{Für wen dieses Buch ist}

Für Entwickler, die EF Core bereits sicher einsetzen: \code{Db\allowbreak{}Context},
Migrationen, LINQ-Abfragen und Dependency Injection werden vorausgesetzt. Kenntnisse der
Interna von EF sind nicht nötig; das Buch führt sie dort ein, wo sie gebraucht werden.
Erfahrung mit PostgreSQL hilft beim Beispiel, aber jedes SQL-Konstrukt wird beim ersten
Auftreten erklärt. Kapitel~\ref{ch:ef-core-s-architecture-from} endet mit einem
vollständigen Mini-Plugin, das alle Bausteine auf einer Seite zeigt, bevor das große
Beispiel beginnt.

\section{Was ein Plugin der Stufe 3 leisten kann}

Ein vollständiges Plugin kann EF Core an praktisch jeder Stelle erweitern: Es kann neue
CLR-Typen als Spaltentypen verfügbar machen, eigene Funktionen in LINQ nach SQL übersetzen,
das Modell automatisch per Convention anpassen, Metadaten über die Fluent API und Attribute
annehmen, das Modell validieren, beim Speichern und beim Öffnen von Verbindungen eingreifen,
eigene Ereignisse protokollieren, Migrationen um eigenes SQL ergänzen und die Werkzeuge von
\code{dotnet ef} erweitern. Das Beispiel in diesem Buch nutzt jede dieser Möglichkeiten.

\begin{tip}[Vorbilder im Ökosystem]
  Bevor man selbst schreibt, lohnt der Blick in bestehende Plugins.
  \textbf{EFCore.NamingConventions} ist ein kleines, sauberes Convention-Plugin. Das
  \textbf{NetTopologySuite}-Plugin im EF-Core-Repository zeigt Type Mappings und
  Übersetzungen. Das \textbf{NodaTime}-Plugin von Npgsql zeigt ein provider-spezifisches
  Plugin, das sich über den Options-Builder des Providers aktivieren lässt.
\end{tip}

% ------------------------------------------------------------
\chapter{Die Architektur von EF Core aus Plugin-Sicht}\label{ch:ef-core-s-architecture-from}

\section{Zwei Container: der interne Service Provider}

Jeder \code{Db\allowbreak{}Context} löst seine Infrastruktur aus einem \textbf{internen}
\code{IService\allowbreak{}Provider} auf. Er ist vom DI-Container der Anwendung getrennt und
enthält mehrere hundert Dienste: Type Mapping, Model Building, Query-Pipeline, Change
Tracking, Update-Pipeline, Migrations. Ein Plugin klinkt sich ein, indem es in diesen
Container zusätzliche Dienste registriert oder vorhandene ersetzt.

\begin{flowdiagram}
  \flowstep{DbContextOptions}{Provider-Extension (UseNpgsql) und Plugin-Extensions
  (UseAcmeTenancy)}
  \flowstep{ServiceProviderCache}{Sucht per Hash einen passenden internen Provider oder baut
  einen}
  \flowstep{ApplyServices()}{Jede Extension registriert ihre Dienste im neuen Provider}
  \flowstep{Interner Service Provider}{Gecacht und von allen passenden Contexts geteilt}
  \flowstep{Model (gecacht)}{Einmal pro Provider gebaut: Conventions, Type Mappings,
  Validierung}
  \flowstep{Query, SaveChanges, Migrations}{Laufen mit den Diensten dieses Providers}
\end{flowdiagram}

Zwei Caches sind dabei entscheidend. Der \textbf{Service Provider} wird anhand der Options
gecacht: Contexts mit gleicher Konfiguration teilen sich einen Provider. Das \textbf{Modell}
wird pro Provider (und Context-Typ) gecacht. Beides macht EF Core schnell, hat aber direkte
Folgen für Plugins, die Kapitel~\ref{ch:the-options-extension} ausführlich behandelt.

\section{Lebensdauern}

Jeder EF-Dienst hat eine feste Lebensdauer. \textbf{Singleton} heißt hier: einmal pro
internem Service Provider, nicht pro Anwendung. \textbf{Scoped} heißt: einmal pro
Context-Instanz. Ein Singleton darf keinen Scoped-Dienst halten, und ein Dienst, der pro
Context unterschiedliche Daten braucht, darf kein Singleton sein. Anhang~\ref{app:A} listet
die Lebensdauern aller Erweiterungspunkte.

\section{Das Dependencies-Muster}

EF-Dienste bekommen ihre Abhängigkeiten nicht einzeln, sondern als ein Parameterobjekt, etwa
\code{Migrations\allowbreak{}Sql\allowbreak{}Generator\allowbreak{}Dependencies} oder
\code{Provider\allowbreak{}Convention\allowbreak{}Set\allowbreak{}Builder\allowbreak{}Dependencies}.
Wer von einem EF-Dienst erbt, übernimmt dieses Objekt im Konstruktor und reicht es weiter.
Neue Abhängigkeiten fügt EF Core dem Objekt hinzu, ohne Konstruktoren zu brechen. Für
Plugins ist das ein Segen: Abgeleitete Dienste bleiben über Minor-Versionen kompatibel.

\begin{sourcecode}[Migrations/\allowbreak{}Internal/\allowbreak{}Tenancy\allowbreak{}Migrations\allowbreak{}Sql\allowbreak{}Generator.\allowbreak{}cs (\allowbreak{}Kopf)]{csharp}
public class TenancyMigrationsSqlGenerator : NpgsqlMigrationsSqlGenerator
{
    public TenancyMigrationsSqlGenerator(
        MigrationsSqlGeneratorDependencies dependencies,     // Parameterobjekt von EF
        INpgsqlSingletonOptions npgsqlSingletonOptions)      // zusätzlich vom Provider
        : base(dependencies, npgsqlSingletonOptions)
    {
    }
\end{sourcecode}

\begin{note}[Code-Stil im Beispiel]
  Alle Klassen im Beispiel verwenden klassische Konstruktoren mit \code{private
  readonly}-Feldern und
  \code{Argument\allowbreak{}Null\allowbreak{}Exception.\allowbreak{}Throw\allowbreak{}If\allowbreak{}Null},
  keine Primary Constructors. Deren Parameter sind im ganzen Klassenkörper veränderbar,
  bieten keinen natürlichen Ort für Validierung und verbergen, welcher Zustand gespeichert
  wird. Klassische Konstruktoren entsprechen außerdem dem Stil von EF Core und Npgsql, von
  deren Klassen ein Plugin erbt. Einzige Ausnahme sind positionelle Records wie
  \code{Rls\allowbreak{}Policy}: Dort erzeugt die Syntax echte \code{init}-Properties.
\end{note}

\section{Erweitern oder ersetzen?}

Für Erweiterungen gibt es zwei grundverschiedene Wege. \textbf{Plugin-Interfaces} wie
\code{IMethod\allowbreak{}Call\allowbreak{}Translator\allowbreak{}Plugin} oder
\code{IConvention\allowbreak{}Set\allowbreak{}Plugin} sind ausdrücklich dafür gedacht. Sie
sind mehrfach registrierbar, EF fragt alle Registrierungen ab, und mehrere Plugins vertragen
sich. \textbf{Ersetzen} bedeutet, einen Einzeldienst wie
\code{IMigrations\allowbreak{}Sql\allowbreak{}Generator} durch eine abgeleitete Klasse
auszutauschen. Das ist manchmal unvermeidlich, aber es kann pro Context immer nur einer
gewinnen.

\begin{important}[Plugin-Interfaces vor Ersetzen]
  Immer zuerst prüfen, ob es ein Plugin-Interface gibt. Wer einen Einzeldienst ersetzt,
  kollidiert mit jedem anderen Plugin, das denselben Dienst ersetzt, und muss von der
  Provider-Implementierung erben. Das bindet das Plugin an einen Provider und an dessen
  interne API. Kapitel~\ref{ch:registering-services} zeigt, wie man trotzdem kombinierbar
  bleibt.
\end{important}

\section{Zum Aufwärmen: ein vollständiges Plugin in fünf Teilen}

Vor dem großen Beispiel kommt hier das kleinste Plugin, das trotzdem vollständig ist. Es
macht PostgreSQLs \code{reverse(\allowbreak{})} in LINQ als
\code{EF.\allowbreak{}Functions.\allowbreak{}Reverse(\allowbreak{}.\allowbreak{}.\allowbreak{}.\allowbreak{})}
verfügbar. Jedes spätere Kapitel vertieft einen seiner fünf Teile:

\begin{sourcecode}[tests/\allowbreak{}Acme.\allowbreak{}Entity\allowbreak{}Framework\allowbreak{}Core.\allowbreak{}Tenancy.\allowbreak{}Tests/\allowbreak{}Minimal\allowbreak{}Plugin\allowbreak{}Tests.\allowbreak{}cs (\allowbreak{}1/\allowbreak{}2,\allowbreak{} gekürzt)]{csharp}
// 1. Die Funktion, die Anwender in LINQ aufrufen; sie wird nie ausgeführt
public static class ReverseDbFunctionsExtensions
{
    public static string Reverse(this DbFunctions _, string value)
        => throw new InvalidOperationException("Reverse can only be used in LINQ queries.");
}

// 2. Der Translator macht aus dem Methodenaufruf eine SQL-Funktion
//    (Kapitel „Methoden und Member übersetzen“)
public class ReverseTranslator : IMethodCallTranslator
{
    // Konstruktor speichert die ISqlExpressionFactory in _sql

    public SqlExpression? Translate(
        SqlExpression? instance, MethodInfo method, IReadOnlyList<SqlExpression> arguments,
        IDiagnosticsLogger<DbLoggerCategory.Query> logger)
        => method == ReverseMethod
            ? _sql.Function("reverse", [arguments[1]], nullable: true,
                argumentsPropagateNullability: [true], typeof(string),
                arguments[1].TypeMapping)
            : null;
}

// 3. Das Plugin übergibt den Translator an EF (Kapitel „Methoden und Member übersetzen“)
public class ReverseTranslatorPlugin : IMethodCallTranslatorPlugin
{
    public ReverseTranslatorPlugin(ISqlExpressionFactory sql)
    {
        Translators = [new ReverseTranslator(sql)];
    }

    public IEnumerable<IMethodCallTranslator> Translators { get; }
}
\end{sourcecode}

\begin{sourcecode}[tests/\allowbreak{}Acme.\allowbreak{}Entity\allowbreak{}Framework\allowbreak{}Core.\allowbreak{}Tenancy.\allowbreak{}Tests/\allowbreak{}Minimal\allowbreak{}Plugin\allowbreak{}Tests.\allowbreak{}cs (\allowbreak{}2/\allowbreak{}2,\allowbreak{} gekürzt)]{csharp}
// 4. Die Options-Extension registriert das Plugin im internen Provider
//    (Kapitel „Die Options-Extension“ und „Dienste registrieren“)
public class ReverseOptionsExtension : IDbContextOptionsExtension
{
    private DbContextOptionsExtensionInfo? _info;

    public DbContextOptionsExtensionInfo Info => _info ??= new ExtensionInfo(this);

    public void ApplyServices(IServiceCollection services)
        => new EntityFrameworkRelationalServicesBuilder(services)
            .TryAdd<IMethodCallTranslatorPlugin, ReverseTranslatorPlugin>();

    public void Validate(IDbContextOptions options)
    {
    }

    // ExtensionInfo: kein Provider, keine Settings -> Hash 0, immer derselbe Provider
}

// 5. Die Use-Methode aktiviert das Plugin
//    (Kapitel „Öffentliche API und Integration in die Anwendung“)
public static DbContextOptionsBuilder UseReverse(this DbContextOptionsBuilder builder)
{
    ((IDbContextOptionsBuilderInfrastructure)builder).AddOrUpdateExtension(
        builder.Options.FindExtension<ReverseOptionsExtension>()
        ?? new ReverseOptionsExtension());
    return builder;
}
\end{sourcecode}

\begin{sourcecode}[Anwendung]{csharp}
options.UseNpgsql(connectionString).UseReverse();

db.Customers.Where(c => EF.Functions.Reverse(c.Name) == "abc");
// SQL: WHERE reverse(c."Name") = 'abc'
\end{sourcecode}

Mehr braucht ein funktionierendes Plugin nicht. EF findet den Translator, weil die Extension
ihn im internen Service Provider registriert, und alle Contexts, die mit
\code{Use\allowbreak{}Reverse(\allowbreak{})} konfiguriert sind, teilen sich diesen
Provider. Weil das Argument auf eine Spalte verweist, braucht die Funktion keinen
Evaluatable-Expression-Filter (Kapitel~\ref{ch:preventing-client-side-evaluation} erklärt,
wann sie einen braucht). Der Test \code{Minimal\allowbreak{}Plugin\allowbreak{}Tests} prüft
die Übersetzung und dass dieselbe Abfrage ohne das Plugin fehlschlägt. Das Mandanten-Plugin
besteht aus genau diesen Teilen, nur aus sehr viel mehr davon.

% ------------------------------------------------------------
\chapter{Das Beispielszenario}\label{ch:the-sample-scenario}

Das Beispiel-Plugin \textbf{Acme.EntityFrameworkCore.Tenancy} macht EF Core mit PostgreSQL
mandantenfähig. Die Trennung der Mandanten erzwingt die Datenbank selbst über \textbf{Row
Level Security} (RLS): Jede Tabelle bekommt eine Policy, die nur Zeilen des Mandanten zeigt,
dessen ID in der aktuellen Datenbank-Session gesetzt ist. Selbst eine fehlerhafte Abfrage
oder rohes SQL kann so keine fremden Daten lesen.

Aus Sicht des Anwenders sieht das so aus:

\begin{sourcecode}[Anwendersicht]{csharp}
public class Order : ITenantOwned               // Marker aus dem Kernpaket
{
    public int Id { get; set; }
    public TenantId TenantId { get; set; }      // stark typisiert, automatisch uuid
    public decimal Total { get; set; }
}

services.AddAcmeTenancy<HttpTenantProvider>();
services.AddDbContext<ShopContext>((sp, o) => o
    .UseNpgsql(connectionString)
    .UseAcmeTenancy(sp, tenancy => tenancy.ApplyQueryFilter()));

// Abfragen und Speichern ohne jede Mandantenlogik im Anwendungscode
db.Orders.Add(new Order { Total = 42 });         // TenantId setzt das Plugin
var mine = await db.Orders.ToListAsync();        // RLS filtert in der Datenbank
\end{sourcecode}

Damit das funktioniert, nutzt das Plugin jeden Erweiterungspunkt dieses Buchs:

\begin{datatable}{H{1.20}H{1.50}H{0.30}}
  \textbf{Anforderung} & \textbf{Umsetzung} & \textbf{Kap.} \\
  \tablehead
  Konfiguration über \code{Use\allowbreak{}Acme\allowbreak{}Tenancy(\allowbreak{})} &
  Options-Extension, Settings-Record, Options-Builder & 4–6 \\
  \tablerule
  \code{Tenant\allowbreak{}Id} als Spaltentyp ohne Konfiguration & Type-Mapping-Plugin,
  Value Converter, JSON-Reader/Writer & 10 \\
  \tablerule
  \code{EF.\allowbreak{}Functions.\allowbreak{}Current\allowbreak{}Tenant(\allowbreak{})} in
  LINQ & Method-Translator, Evaluatable-Expression-Filter & 12–13 \\
  \tablerule
  \code{tenant\allowbreak{}Id.\allowbreak{}Is\allowbreak{}Empty} in LINQ & Member-Translator
  & 12 \\
  \tablerule
  RLS für alle \code{ITenant\allowbreak{}Owned}-Entitäten & Convention mit Index und
  optionalem Query-Filter & 8 \\
  \tablerule
  Fluent API, Attribut \code{[No\allowbreak{}Tenant\allowbreak{}Isolation]} & Annotations,
  Attribut-Convention, ConfigurationSource & 7–8 \\
  \tablerule
  Fehlkonfiguration früh erkennen & Validierungs-Convention, eigene Warnungen & 9, 16 \\
  \tablerule
  TenantId automatisch setzen & SaveChanges-Interceptor & 15 \\
  \tablerule
  Mandant in der DB-Session & Connection-Interceptor & 15 \\
  \tablerule
  \code{ENABLE ROW LEVEL SECURITY}, \code{CREATE POLICY} & Annotation-Provider,
  Migrations-SQL-Generator & 17–18 \\
  \tablerule
  Rollen und Rechte in Migrationen & Eigene Migrations-Operationen & 19 \\
  \tablerule
  Lesbare Snapshots & Annotation-Code-Generator, Design-Time-Services & 20 \\
  \tablerule
  Compiled Models und AOT & \code{Default}-Instanzen, Source-Generatoren & 21 \\
  \tablerule
  \code{EF.\allowbreak{}Functions.\allowbreak{}Is\allowbreak{}Not\allowbreak{}Distinct\allowbreak{}From},
  \code{EF.\allowbreak{}Functions.\allowbreak{}Every} & Eigener SQL-Knoten,
  Aggregat-Translator & 37–38 \\
  \tablerule
  Mandanten-Property mit beliebigem Namen &
  \code{Is\allowbreak{}Tenant\allowbreak{}Property(\allowbreak{})}, Laufzeit-Annotation & 39
  \\
  \tablerule
  RLS aus bestehender Datenbank scaffolden & Eigene Database-Model-Factory & 40 \\
  \tablerule
  SQL Server & Zweites Provider-Paket auf demselben Kern & 41–42 \\
\end{datatable}

\section{Paketarchitektur}

Das Plugin besteht aus drei NuGet-Paketen. Das \textbf{Kernpaket} enthält alles, was
unabhängig von der Datenbank ist und was auch Assemblies mit Entitätsklassen referenzieren
dürfen, ohne einen Provider mitzuziehen. Das \textbf{PostgreSQL-Paket} ist die vollständige
Implementierung, an der dieses Buch alles erklärt. Das \textbf{SQL-Server-Paket} nutzt
denselben Kern und zeigt in Kapitel~\ref{ch:a-second-provider-sql-server}, was ein zweiter
Provider selbst umsetzen muss.

\begin{sourcecode}[Struktur der drei Pakete]{text}
Acme.EntityFrameworkCore.Tenancy               Kern (EF Core relational, ohne Provider)
├─ TenantId, TenantIdConverter, TenantIdJsonReaderWriter, TenancyException
├─ ITenantOwned, ITenantProvider, NoTenantIsolationAttribute
├─ EF.Functions.CurrentTenant(), TenancyDbFunctionsAttribute
├─ Metadata/         Laufzeit-Annotation, IsTenantProperty(), Kandidatensuche
├─ Conventions/      Value Generation für TenantId-Schlüssel
├─ Query/            Evaluatable-Expression-Filter (für alle Pakete)
├─ Interceptors/     SaveChanges-Interceptor, abstrakter Session-Interceptor
├─ Diagnostics/      Event-IDs, Logging-Definitionen, Logger-Extensions
└─ DependencyInjection/  AddAcmeTenancy<T>()

Acme.EntityFrameworkCore.Tenancy.PostgreSQL    vollständige Implementierung
├─ Infrastructure/   Options-Extension, Builder, Singleton-Options
├─ Storage/          TenantIdTypeMapping, Type-Mapping-Plugin
├─ Query/            Translators, eigener SQL-Knoten, Aggregat, SQL-Generator
├─ Metadata/         RLS-Annotations, RlsPolicy, Annotation-Provider
├─ Builders/         Fluent API (HasRowLevelSecurity, HasPolicy, ...)
├─ Conventions/      Isolation, Attribut, Validierung, Convention-Set-Plugin
├─ Migrations/       SQL-Generator, eigene Operationen
├─ Design/           Code-Generatoren, Design-Time-Services
├─ Scaffolding/      RLS und Policies aus der Datenbank lesen
└─ Runtime/          TenantConnectionInterceptor

Acme.EntityFrameworkCore.Tenancy.SqlServer     zweiter Provider
                                               (Kapitel „Ein zweiter Provider: SQL Server“)
└─ Infrastructure/, Storage/, Query/, Metadata/, Conventions/, Migrations/, Runtime/
\end{sourcecode}

\begin{tip}
  Die Ordnerstruktur folgt der von EF Core selbst (Infrastructure, Storage, Query, Metadata,
  Migrations, Design). Wer EF Core kennt, findet sich sofort zurecht, und beim Nachschlagen
  im EF-Quellcode weiß man, wo das Gegenstück liegt.
\end{tip}
